\subsection{Groundwater PFAS transport and integration}
\label{sec:pfas}

PFAS groundwater transport is simulated with a MODFLOW~6 groundwater-transport (GWT) model run on the
calibrated flow field through the GWF6--GWT6 exchange, so the transport uses the exact flow solution.
The model solves advection (total-variation-diminishing), dispersion ($\alpha_L=10$~m, $\alpha_T=1$~m
at the 250~m grid), and \emph{Freundlich} sorption (MODFLOW~6 MST package; $K_f=0.005$, exponent
$n=0.8$), the standard PFAS isotherm, giving a concentration-dependent retardation factor. This first
implementation simulates PFOS groundwater transport; the multi-analyte source fingerprinting reported in
Section~\ref{sec:results} is an observation-based decomposition independent of the transport model, and
compound-specific groundwater transport with per-analyte sorption isotherms is a priority for future work.

\paragraph{Prescribed source.} The PFAS source is set from measurements rather than fitted: there is
no calibrated groundwater concentration multiplier. Constant-concentration cells are placed at the
source-zone grid cells whose measured groundwater PFOS exceeds $10^4$~ng~L\textsuperscript{-1}, at
their observed concentration. The single peak sample reaches $1.5\times10^6$~ng~L\textsuperscript{-1};
the source is capped at $10^5$~ng~L\textsuperscript{-1}, because that one peak makes the
six-order-of-magnitude front numerically intractable and because the log-space validation tests the
plume \emph{shape} rather than the absolute peak (the source magnitude carries the single-sample
uncertainty). A small ambient background ($10$~ng~L\textsuperscript{-1}, near the observed median)
both removes the Freundlich $C^{n-1}$ singularity at $C\to0$ and compresses the dynamic range;
adaptive time stepping carries the stiff transient over a 40-year horizon (the tannery legacy). The
transport runs as a transient on the calibrated \emph{steady} flow field, a simplification justified
by the multi-decadal source history being long relative to seasonal flow variation but which omits
non-stationary pumping and recharge over that period --- a limitation for the transient plume timing,
though the validation compares time-aggregated plume concentrations. The field is thus a constrained
prediction: bounded by the prescribed source, the cap, and the background, but with no fitted
concentration term. The groundwater source magnitude and location are therefore fixed by measurement;
the only quantity fitted to in-stream observations is the fraction of the discharged load that reaches
the stream --- the joint-calibration scalar $g$ (Section~\ref{sec:results}) --- not the source itself.

\paragraph{Validation.} The modeled groundwater plume is validated against the $846$ measured
groundwater PFOS observations in the watershed --- assigned to grid cells by our extraction pipeline,
which snaps the inventory to the MODFLOW~6 grid built in Section~\ref{sec:pfas} --- excluding the
source-zone cells (Table~\ref{tab:validation}).
Skill is reported in log space (root-mean-square error, bias, fraction within a factor of ten), the
appropriate metric for concentrations spanning many orders of magnitude, and is \emph{stratified by
concentration band}: a factor-of-ten miss near the $\sim$2~ng~L\textsuperscript{-1} detection limit is
measurement noise, whereas the same miss at tens of ng~L\textsuperscript{-1} is a model error, and a
single pooled statistic conflates the two.

\paragraph{Porewater discharge validation.} A second, more direct test validates the groundwater
contribution at the discharge interface itself, rather than only in the far-field plume: streambed
\emph{pore water} sampled along the lower mainstem (EGLE North Kent inventory; $43$ sites, all within
$100$~m of a channel and concentrated on the source-bearing reaches) measures the PFAS concentration of
groundwater as it enters the stream. The modeled aquifer concentration at each pore-water cell is
compared to the measurement in log space by concentration band. This is the most direct available
check of the modeled pathway, since pore water samples the discharging groundwater in place rather than
the plume upgradient of it.

\paragraph{Integration with the surface-water pathway.} The groundwater-discharged PFAS reaches the
stream through the MODFLOW~6 streamflow-transport (SFT) package, which routes it downstream along the
SFR channel network. The groundwater PFAS \emph{load} delivered to each reach is the reach gaining
flux (the SFR groundwater-exchange budget term, m\textsuperscript{3}~d\textsuperscript{-1}) times the
aquifer PFAS concentration at that cell (ng~L\textsuperscript{-1}), excluding the artificial
constant-concentration source cells. Summed over the network and combined with the
SWAT\textsuperscript{+} surface streamflow, this groundwater contribution to in-stream PFOS adds to
the land-derived load from the surface-water engine, closing the continuous
surface-water-plus-groundwater PFAS mass balance.
